\section{概念深挖：把 Lecture 06 变成可验证研究命题}

到这里，课程观点已经被转成训练和部署流程；最后一步是把“经验上如此”改写成可证伪的研究命题。本节会明确目标函数的差异、通信在不同博弈中的作用域，以及每条主张所需的最小对照实验，使读者能够据此设计复现实验而不是只复述结论。

\subsection{Minimax 与 Population Best Response 的形式化差异}
为了避免概念漂移，可以把两者写成不同优化问题。
Minimax（以二人零和为例）通常写作：
\[
\pi^\star = \arg\max_{\pi_i}\min_{\pi_{-i}} \mathbb{E}_{a\sim(\pi_i,\pi_{-i})}[u_i(a)].
\]
Population Best Response 则是：
\[
\pi^\star_{\text{PBR}}=\arg\max_{\pi_i}\mathbb{E}_{\pi_{-i}\sim P_{\text{pop}}}\left[\mathbb{E}_{a\sim(\pi_i,\pi_{-i})}[u_i(a)]\right].
\]
两式的差异不在符号，而在 ``谁定义了对手分布''。Minimax 的对手是最坏情形，PBR 的对手是群体分布。多人协作任务里，后者往往才是部署目标。

\begin{knowledgebox}{课程中的隐含逻辑}
Brown 并非否定 Minimax，而是强调其适用域。离开二人零和后，继续把 Minimax 当总目标，等价于把错误 inductive bias 注入训练过程。
\end{knowledgebox}

\subsection{为什么 ``二人零和中 cheap talk 无效'' 值得重视}
课程给出的是一个非常工程化的证明直觉：若通信动作对发送方有利，接收方会忽略；若通信动作对发送方有害，发送方不会说；唯一稳定情况是通信无效。
这个结论可转化为系统设计原则：在严格对抗任务中，花大量预算做自然语言协商机制，边际收益通常很低。

\begin{importantbox}{设计原则}
当任务接近二人零和时，优先投入 ``策略鲁棒性与对抗评测''，而非 ``多轮语言协商复杂度''。
\end{importantbox}

反过来，在非零和或协调博弈中，通信可显著改变可达均衡集合。此时 ``是否能通信'' 不是关键，``通信后是否能稳定收敛'' 才是关键。

\subsection{从课堂陈述到可检验实验：命题矩阵}
为了让讲义可直接用于研究立项，下面把课程核心观点转成实验命题：

\begin{longtable}{p{3.0cm}p{4.6cm}p{4.9cm}}
\toprule
\textbf{命题} & \textbf{最小实验设计} & \textbf{可接受结论标准} \\
\midrule
P1：自博弈外推受限 & 在同任务中比较 self-play-only 与 human-augmented 两条线，并做跨群体评测 & self-play-only 在跨群体指标显著退化 \\
P2：PPO 在不完美信息混合策略上不稳 & 用同一环境对比 PPO 与 regret 系，统计 exploitability 下降曲线 & regret 系在稳定性和收敛速率上占优 \\
P3：Consensus 只在短答案优势明显 & 分别在短答案与长文生成任务中测 consensus 增益 & 短答案提升显著，长文本提升有限或回撤 \\
P4：Best-of-N 依赖验证器质量 & 固定采样预算，替换不同 verifier 并测最终正确率 & 验证器精度与系统增益呈强正相关 \\
P5：人群分布决定协作最优策略 & 在不同文化/偏好分组上训练并交叉测试 & 组内最优策略在组外交叉明显失配 \\
\bottomrule
\end{longtable}

\begin{warningbox}{实验汇报中的常见失真}
只报 ``平均分提升'' 而不报跨群体方差，会掩盖部署风险。对多 Agent 系统，方差往往比均值更有决策价值。
\end{warningbox}

\subsection{研究者常问问题（FAQ）}
\textbf{Q1：是否还值得做纯 self-play？}
值得，但要把目标限定在对抗稳健或可控博弈子空间，不要默认它会自动学到人类协作规范。

\textbf{Q2：多 Agent 与工具调用有什么关系？}
工具调用可视为 ``异构专家协作'' 的最简形式，通常比开放式自由协商更易控、也更容易评测。

\textbf{Q3：如何判断是否该加更多 agent？}
看边际收益曲线。如果新增 agent 只带来很小质量提升但显著增加时延与成本，应回退为更简单架构。

\textbf{Q4：课程中提到的 ``现在是好时机'' 如何理解？}
因为通信门槛已由 LLM 大幅降低，真正瓶颈转移到协议可靠性和群体泛化，这是可做出新贡献的空白地带。

\subsection{本章小结}
Lecture 06 的最大价值是把讨论从 ``多 Agent 要不要做'' 推进到 ``在什么目标和分布下做才对''。当命题可检验、指标可复现时，这门课就从观点变成了研究工具。

